\chapter{Wellenoptik}

\section{2024-06-06 - Interferenz Gitter Versuch}

\subsection{Beobachtung}

Abstand zum Schirm: 27cm Abstand der Maxima: 12cm

\subsection[title={Auswertung},reference={auswertung}]

\startblockquote
Algemein sind folgende Formeln bekannt:
\stopblockquote

\startformula 
\sin{\alpha} = \frac{\lambda}{g} \quad \text{und} \quad \tan{\alpha}=\frac{a}{l}
\stopformula

Wobei $\lambda$ die Wellenlaenge ist.

\startblockquote
Gitter: 500 Spalten pro Millimeter
\stopblockquote

\startformula 
g=\frac{1\cdot 10^{-3}m}{500}=2\cdot 10^{-6}m
 \stopformula

\startitemize[packed]
\item
  $2a_1=0,12m; \quad a_1 = 0,06m; \quad l=27cm=0,27m$
\stopitemize

\startformula 
\startalign
\lambda \NC =g \cdot \NC \sin{(\tan^{-1}{(\frac{a}{l})})}\NR
 \NC = (2\cdot 10^{-6}) \cdot \NC \sin{(\tan^{-1}{(\frac{0,12}{0,27})})}\NR
 \NC =434\cdot 10^{-9}m
\stopalign
 \stopformula





\subsubsection[title={Versuch
Wiederholung},reference={versuch-wiederholung}]

\startformula 
\startalign
2a_2 = 0.127m; \quad a_2 = 0.635m; \quad l = 0.38m \\
\stopalign
 \stopformula

Berechnung der Wellenlaenge $\lambda$:

\startformula 
\startalign
\lambda \NC =g \cdot \NC \sin{(\tan^{-1}{(\frac{a}{l})})}\NR
 \NC = (2\cdot 10^{-6}) \cdot \NC \sin{(\tan^{-1}{(\frac{0,07}{0,38})})}\NR
 \NC = 6,34\cdot 10^{-7}m=634nm \NR
\stopalign
 \stopformula

\startframedtask
$2a$ ist zwischen den Maxima der Ordnung $n$. Also von einem Maxima bis
zur mitte ist nur $a$
\stopframedtask

\subsubsection[title={Zweite Runde},reference={zweite-runde}]

\startitemize[packed]
\item
  2024-06-18
\stopitemize

\subsubsection[title={Messung der verschiedenen Wellen /
LED's},reference={messung-der-verschiedenen-wellen-leds}]

\starttable[|l|l|l|l|]
\NC LED \VL Wellenlaenge in nm \VL Abstand 1. Ordnung in cm \VL A. 2. Ordnung \SR
\HL
\NC Rot \VL 632 \VL 10,3 \VL - \SR
\NC Grün \VL 514 \VL 8,5 \VL 18.8 \SR
\NC Blau \VL 463 \VL 7,5 \VL 15,7 \SR
\stoptable

\startformula 
g=\frac{1\cdot 10^{-3}m}{500}=2\cdot 10^{-6}m
 \stopformula

\subsubsection[title={Rot},reference={rot}]

\subsubsection[title={1. Ordnung},reference={ordnung}]

\startformula 
2a = 0.103m; \quad a = 0.0515m; \quad l = 0.15m
 \stopformula

Berechnung der Wellenlaenge $\lambda$:

\startformula 
\startalign
\lambda \NC =\frac{g}{n} \cdot \sin{(\tan^{-1}{(\frac{a_n}{l})})}\NR
\NC = (2\cdot 10^{-6}) \cdot \sin{(\tan^{-1}{(\frac{0,0515}{0,15})})}\NR
\NC = 6,49\cdot 10^{-7}m\NR
\stopalign
 \stopformula







